%% Abstract
%%
%% Traduction anglaise fidèle et de qualité du résumé de la recherche écrit en français et non une traduction littérale.
%%

\chapter*{ABSTRACT}\thispagestyle{headings}
\addcontentsline{toc}{compteur}{ABSTRACT}

\begin{otherlanguage}{english}

Three-dimensional multilayer ice accretion simulations on aircraft present significant numerical challenges because they require the successive coupling of aerodynamic, droplet-impingement, thermodynamic, and geometry-evolution models while repeatedly updating an increasingly complex iced geometry. In particular, explicit surface-displacement techniques and repeated body-fitted remeshing can become sensitive to mesh quality as complex 3D ice features develop. The objective of this thesis is to improve the robustness of 3D multilayer ice accretion simulations through the development and assessment of complementary numerical methodologies within \textsc{CHAMPS}, an in-house multiphysics software being developed at Polytechnique Montréal.

First, a Lagrangian particle-tracking methodology is developed for droplet-impingement prediction. The approach combines implicit trajectory integration, accelerated intersection detection, adaptive droplet seeding, and surface interpolation techniques to improve computational efficiency and the prediction of the amount of water collected at the surface. Supercooled Large Droplet (SLD) effects, including droplet deformation, splashing, and secondary-droplet tracking, are incorporated into both Eulerian and Lagrangian formulations. The resulting predictions show good agreement with reference data and demonstrate consistency between the Eulerian and Lagrangian approaches. The comparison also highlights their complementary characteristics: the Eulerian formulation provides more predictable computational cost and efficient parallelization for large 3D problems, whereas the Lagrangian formulation provides direct access to individual droplet trajectories and is particularly attractive for incorporating detailed droplet-scale physics.

Second, a continuous level-set methodology is developed for 3D multilayer ice accretion. The ice--air interface is represented implicitly as the zero contour of a signed-distance field defined on an auxiliary volume mesh. The methodology employs a semi-implicit time-integration scheme (Rosenbrock-Wanner) to advect the interface using a velocity determined from the icing quantities, and subsequently uses the MMG library to extract a conformal surface representing the resulting geometry. Benchmark calculations demonstrate the accuracy of the solver and its ability to reproduce complex 3D ice structures, including experimentally observed scallops, while providing increased geometric robustness compared with explicit surface-node displacement. This advantage becomes particularly important as geometric complexity increases, since the implicit representation avoids surface-mesh entanglement and reduces the need for heuristic geometric corrections during successive accretion layers.

Finally, the advances made in droplet tracking and level-set geometry evolution are integrated into an icing framework based on a hybrid immersed-boundary approach, in which the aerodynamic and droplet-impingement solutions are decoupled from a body-fitted volume mesh conforming to the evolving geometry. Thermodynamic calculations and geometry evolution, however, retain a conformal surface representation to preserve the advantages of the level-set geometry-evolution methodology. Verification and ice-accretion calculations demonstrate encouraging agreement between body-fitted and immersed-boundary predictions for droplet impingement and ice growth. The results also highlight the importance of local mesh resolution and near-wall aerodynamic and thermal treatments. In particular, adaptive mesh refinement is identified as an important component for efficiently extending the immersed-boundary methodology to complex 3D configurations.

Overall, this thesis establishes a set of complementary numerical tools for improving the robustness and flexibility of 3D ice accretion simulations. The results indicate that level-set geometry evolution provides a robust foundation for multilayer simulations, while Lagrangian droplet tracking and immersed-boundary methods provide additional capabilities for detailed droplet modeling and the treatment of complex evolving geometries, respectively.

\end{otherlanguage}